\clearpage
\SourceLeaf{228}\pdfbookmark[1]{Zweiter Abschnitt.}{liber-secundus-sectio-secunda}
\begin{center}
{\large ZWEITER ABSCHNITT.}\par\medskip
\emph{Bestimmung der Bahn aus vier Beobachtungen, von denen nur zwei vollständig sind.}\par\bigskip
\hypertarget{native-228-article164}{}\pdfbookmark[2]{164.}{article-164}164.
\end{center}
Schon zu Beginn des zweiten Buches (Artikel 115) haben wir erklärt, dass der Anwendungsbereich der im vorhergehenden Abschnitt behandelten Aufgabe auf Bahnen beschränkt ist, deren Neigung weder verschwindet noch zu klein ist, und dass die Bestimmung schwach geneigter Bahnen auf vier Beobachtungen beruhen muss. Vier vollständige Beobachtungen entsprechen jedoch acht Gleichungen, und es gibt nur sechs Unbekannte. Die Aufgabe wäre somit überbestimmt. Daher muss man bei zwei Beobachtungen die Breiten (oder Deklinationen) beiseitelassen, damit die übrigen Angaben genau erfüllt werden können. So ergibt sich die Aufgabe, der dieser Abschnitt gewidmet ist. Die hier angegebene Lösung lässt sich auf schwach geneigte Bahnen ebenso erfolgreich anwenden wie auf Bahnen mit beliebig großer Neigung. Auch hier muss man, wie bei der Aufgabe im vorhergehenden Abschnitt, den Fall, in dem die ungefähren Abmessungen der Bahn bereits bekannt sind, von der ersten Bestimmung einer bisher völlig unbekannten Bahn unterscheiden. Wir beginnen mit dem ersteren Fall.
